\documentclass[]{scrartcl}
\usepackage[a4paper,top=0.3in,bottom=0.2in,left=0.3in,right=0.3in]{geometry}
\usepackage{tabulary}
\usepackage{graphicx}
\usepackage{amssymb}
\usepackage{array}
\usepackage{background}
\usepackage[utf8]{inputenc}
\usepackage{eurosym}
\thispagestyle{empty}

\newcommand{\headentry}[1]{\parbox{18.6cm}{#1}}
\newcolumntype{P}[1]{>{\centering\arraybackslash}p{#1}} %uses package 'array'

\begin{document}
	\includegraphics[width=140pt]{Latex_Logo.pdf} \par \smallskip
	\sffamily

	%HEADER
	\vspace{0.2cm}
	\headentry{\huge \textbf{Arbeitszeitdokumentation} \hspace*{\fill} \Large \textbf{Monat / Jahr:} \underline{\parbox{5.0cm}{\centering !month / !year }}} \par \medskip
	\headentry{\Large \textbf{Name des Mitarbeiters/der Mitarbeiterin:} \hspace*{\fill} \underline{\parbox{8.5cm}{!employeeName }}} \par \medskip
	\headentry{\Large \textbf{Personalnummer:} \hspace*{\fill} \underline{\parbox{4.5cm}{!employeeID} \parbox{3.85cm}{\centering !workingArea}}} \par \medskip %This is the KIT style for workingArea => GF: $\Box$ UB: $\boxtimes$
	\headentry{\Large \textbf{Institut/Organisationseinheit:} \hspace*{\fill} \underline{\parbox{8.5cm}{!department }}} \par \medskip
	\headentry{\Large \textbf{Vertraglich vereinbarte Arbeitszeit:} \hspace*{\fill} \parbox{8.5cm}{\underline{\parbox{2.35cm}{\centering !workingTime}} \parbox{6cm}{\centering \raggedleft \textbf{Stundensatz:} \underline{\parbox{2.25cm}{\centering !wage \euro} } } } } \par \medskip

	%BODY
	\large
	\begin{center}
		\begin{tabular}{| P{6.7cm} | P{2cm} | P{1.8cm} | P{1.8cm} | P{1.8cm} | P{2.4cm} |}
			\hline
			%Line 1
			\textbf{T\"atigkeit (Stichwort, Projekt)}
			& \textbf{Datum}
			& \textbf{Beginn}
			& \textbf{Ende}
			& \textbf{Pause}
			& \textbf{Arbeitszeit\textsuperscript{1}}\\
			\hline
			%Line 2
			%empty
			& \textbf{(tt.mm.jj)}
			& \textbf{(hh:mm)}
			& \textbf{(hh:mm)}
			& \textbf{(hh:mm)}
			& \textbf{(hh:mm)}\\
			\hline
			%Line 3
			!action
			& !date
			& !begin
			& !end
			& !break
			& !dayTotal\\
			\hline
			%Line 4
			!action
			& !date
			& !begin
			& !end
			& !break
			& !dayTotal\\
			\hline
			%Line 5
			!action
			& !date
			& !begin
			& !end
			& !break
			& !dayTotal\\
			\hline
			%Line 6
			!action
			& !date
			& !begin
			& !end
			& !break
			& !dayTotal\\
			\hline
			%Line 7
			!action
			& !date
			& !begin
			& !end
			& !break
			& !dayTotal\\
			\hline
			%Line 8
			!action
			& !date
			& !begin
			& !end
			& !break
			& !dayTotal\\
			\hline
			%Line 9
			!action
			& !date
			& !begin
			& !end
			& !break
			& !dayTotal\\
			\hline
			%Line 10
			!action
			& !date
			& !begin
			& !end
			& !break
			& !dayTotal\\
			\hline
			%Line 11
			!action
			& !date
			& !begin
			& !end
			& !break
			& !dayTotal\\
			\hline
			%Line 12
			!action
			& !date
			& !begin
			& !end
			& !break
			& !dayTotal\\
			\hline
			%Line 13
			!action
			& !date
			& !begin
			& !end
			& !break
			& !dayTotal\\
			\hline
			%Line 14
			!action
			& !date
			& !begin
			& !end
			& !break
			& !dayTotal\\
			\hline
			%Line 15
			!action
			& !date
			& !begin
			& !end
			& !break
			& !dayTotal\\
			\hline
			%Line 16
			!action
			& !date
			& !begin
			& !end
			& !break
			& !dayTotal\\
			\hline
			%Line 17
			!action
			& !date
			& !begin
			& !end
			& !break
			& !dayTotal\\
			\hline
			%Line 18
			!action
			& !date
			& !begin
			& !end
			& !break
			& !dayTotal\\
			\hline
			%Line 19
			!action
			& !date
			& !begin
			& !end
			& !break
			& !dayTotal\\
			\hline
			%Line 20
			!action
			& !date
			& !begin
			& !end
			& !break
			& !dayTotal\\
			\hline
			%Line 21
			!action
			& !date
			& !begin
			& !end
			& !break
			& !dayTotal\\
			\hline
			%Line 22
			!action
			& !date
			& !begin
			& !end
			& !break
			& !dayTotal\\
			\hline 
			%Leerzeile
			\multicolumn{6}{c}{\thinspace}\\
			%Urlaub
			\cline{3-6}
			\multicolumn{2}{c}{\thinspace}
			& \multicolumn{3}{|c|}{\centering \textbf{Urlaub anteilig:}}
			& !vacation\\
			\cline{3-6}
			%Summe
			\multicolumn{2}{c}{\thinspace}
			& \multicolumn{3}{|c|}{\centering \textbf{Summe:}}
			& !sum\\
			\cline{3-6}
			%SollArbeitszeit
			\multicolumn{2}{c}{\thinspace}
			& \multicolumn{3}{|c|}{\centering \textbf{monatliche Soll-Arbeitszeit:}}
			& !workingTime\\
			\cline{3-6}
			%Übertrag Vormonat
			\multicolumn{2}{c}{\thinspace}
			& \multicolumn{3}{|c|}{\centering \textbf{Übertrag vom Vormonat:}}
			& !carryPred\\
			\cline{3-6}
			%Übertrag Folgemonat
			\multicolumn{2}{c}{\thinspace}
			& \multicolumn{3}{|c|}{\centering \textbf{Übertrag in den Folgemonat:}}
			& !carrySucc\\
			\cline{3-6}
		\end{tabular}
	\end{center}

	%FOOTER
	\par \bigskip \bigskip \medskip
	\headentry{\large Ich bestätige die Richtigkeit der Angaben: \hspace*{\fill} $\overline{{\parbox{5.75cm}{\normalsize Datum, Beschäftigte/r} } }$ } \par \medskip
	\headentry{\normalsize Nach \textbf{$\S 17$ Mindestlohngesetz (MiLoG)} müssen für geringfügig entlohnte und kurzfristig beschäftigte Arbeitnehmer/innen u.a. Beginn, Ende und Dauer der täglichen Arbeitszeit aufgezeichnet und für Kontrollzwecke mindestens 2 Jahre am Ort der Beschäftigung aufbewahrt werden!} \par \bigskip \bigskip
	\headentry{\hspace*{\fill} geprüft: $\overline{{\parbox{5.75cm}{\normalsize Datum, Dienstvorgesetzte/r} } }$} \par \medskip
	\rule{6cm}{0.2pt} \par \smallskip
	\headentry{\textsuperscript{1} Summe in vollen Stunden und Minuten ohne Pause (Std:Min); bei Abwesenheit können auch folgende Kürzel eingetragen werden: U=Urlaub, K=Krankheit, F=Feiertag, S=Sonstiges}

	%BACKGROUND
	\SetBgContents{K\_IPD\_AZDoku\_01\_01-20}
	\SetBgPosition{-2.4cm, -29.2cm}
	\SetBgColor{black}
	\SetBgOpacity{1.0}
	\SetBgAngle{90.0}
	\SetBgScale{0.8}
\end{document}
